% ============================================================
%  Übungsaufgaben Template (my_dhbw Stil, uebung_*.tex)
%  Header "Übungsblatt", grün eingefärbte <details>-Lösungen.
%  Renderer: pdflatex (zweimal)
% ============================================================
\documentclass[11pt, a4paper]{article}

\usepackage[ngerman]{babel}
\usepackage{fontspec}
\setmainfont[
  Path=/usr/share/fonts/TTF/,
  Extension=.ttf,
  BoldFont=DejaVuSans-Bold,
  ItalicFont=DejaVuSans-Oblique,
  BoldItalicFont=DejaVuSans-BoldOblique
]{DejaVuSans}
\setsansfont[
  Path=/usr/share/fonts/TTF/,
  Extension=.ttf,
  BoldFont=DejaVuSans-Bold,
  ItalicFont=DejaVuSans-Oblique,
  BoldItalicFont=DejaVuSans-BoldOblique
]{DejaVuSans}
\setmonofont[Path=/usr/share/fonts/TTF/,Extension=.ttf]{DejaVuSansMono}
\usepackage[top=2cm, bottom=2cm, left=2.4cm, right=2.4cm, footskip=1cm]{geometry}
\usepackage{amsmath, amssymb}
\usepackage{xcolor}
\usepackage{enumitem}
\usepackage{fancyhdr}
\usepackage{lastpage}
\usepackage{tabularx}
\usepackage{booktabs}
\usepackage{microtype}
\usepackage{parskip}

\usepackage{array}
\usepackage{longtable}
\usepackage{ltablex}
\keepXColumns
\usepackage{colortbl}
\usepackage[most,breakable]{tcolorbox}
\usepackage{listings}
\usepackage[normalem]{ulem}
\usepackage{pdflscape}
\usepackage{titlesec}
\usepackage[colorlinks=true, linkcolor=accent, urlcolor=accent]{hyperref}

\renewcommand{\familydefault}{\sfdefault}

% ── Theme ─────────────────────────────────────────────────────
\definecolor{accent}{RGB}{0, 60, 113}
\definecolor{primaryblue}{RGB}{0, 60, 113}
\definecolor{loesung}{RGB}{0, 100, 0}
\definecolor{codebg}{RGB}{245, 245, 245}
\definecolor{figurebg}{RGB}{232, 241, 255}
\definecolor{tablehintbg}{RGB}{240, 255, 240}
\definecolor{solutionbg}{RGB}{240, 252, 240}
\definecolor{rulegray}{RGB}{180, 180, 180}
\definecolor{tableheadbg}{RGB}{0, 60, 113}

% ── Header/Footer ─────────────────────────────────────────────
\pagestyle{fancy}
\fancyhf{}
\renewcommand{\headrulewidth}{0.4pt}
\fancyhead[L]{\footnotesize\color{accent}\nichttitle}
\fancyhead[R]{\footnotesize\color{accent}\nichtsubtitle}
\fancyfoot[R]{\footnotesize\color{accent}Blatt \thepage\,/\,\pageref{LastPage}}

% ── Typografie ────────────────────────────────────────────────
\titleformat{\section}
    {\color{accent}\large\bfseries}{}{0pt}{}
    [\vspace{-4pt}\color{accent}\rule{\linewidth}{1.5pt}]
\titleformat{\subsection}
    {\normalsize\bfseries\color{accent!85!black}}{}{0pt}{}{}
\titleformat{\subsubsection}
    {\small\bfseries\color{accent!70!black}}{}{0pt}{}{}
\titlespacing*{\section}{0pt}{10pt}{3pt}
\titlespacing*{\subsection}{0pt}{6pt}{2pt}
\titlespacing*{\subsubsection}{0pt}{4pt}{1pt}

\setlength{\parindent}{0pt}
\setlength{\parskip}{6pt}
\renewcommand{\arraystretch}{1.2}
\setlist[itemize]{noitemsep, topsep=3pt, parsep=0pt, leftmargin=1.4em}
\setlist[enumerate]{noitemsep, topsep=3pt, parsep=0pt, leftmargin=1.6em}

% ── Boxen: solutionbox = Lösungsbox in grün ──────────────────
\newtcolorbox{figurebox}[1]{
  colback=figurebg, colframe=accent!50,
  title={Abbildung: #1}, fonttitle=\small\bfseries,
  breakable, before skip=6pt, after skip=6pt
}
\newtcolorbox{tablehintbox}[1]{
  colback=tablehintbg, colframe=green!50!black!50,
  title={Tabelle: #1}, fonttitle=\small\bfseries,
  breakable, before skip=6pt, after skip=6pt
}
\newtcolorbox{solutionbox}[1]{
  enhanced,
  colback=solutionbg, colframe=loesung,
  title={#1}, fonttitle=\small\bfseries,
  coltitle=white, colbacktitle=loesung,
  breakable, before skip=8pt, after skip=8pt
}

\lstset{
  basicstyle=\ttfamily\small,
  backgroundcolor=\color{codebg},
  breaklines=true,
  frame=single, framerule=0pt,
  xleftmargin=4pt, xrightmargin=4pt,
  aboveskip=6pt, belowskip=6pt,
  keepspaces=true
}

\providecommand{\nichttitle}{Übungsblatt}
\providecommand{\nichtsubtitle}{}

\renewcommand{\nichttitle}{Optimierungsverfahren}
\renewcommand{\nichtsubtitle}{Übungsblatt 2}


\begin{document}

\begin{center}
{\Large\bfseries\color{accent}\nichttitle}
\ifx\nichtsubtitle\empty\else
  \\[2pt]{\normalsize\color{accent!80!black}\nichtsubtitle}
\fi
\\[2pt]
\color{accent}\rule{0.4\linewidth}{0.8pt}\color{black}
\end{center}
\vspace{4pt}

% ── BODY ──────────────────────────────────────────────────────

\section{Übungsblatt 2 — 1. Grundlagen}


\noindent\textcolor{rulegray}{\rule{\linewidth}{0.4pt}}


\paragraph{Aufgabe 1 — Problemmodellierung Lieferdrohne \textit{(15 Punkte)}}\mbox{}\\[2pt]

Ein Logistik-Startup will eine Lieferdrohne so konfigurieren, dass die \textbf{Energie pro Lieferung} minimiert wird. Die Drohne hat eine einstellbare \textbf{Flughöhe} $x_1 \in [50, 200]$ m und eine \textbf{Reisegeschwindigkeit} $x_2 \in [5, 20]$ m/s. Die Außentemperatur $T$ und die Windgeschwindigkeit $w$ sind nicht steuerbar. Aus Sicherheitsgründen muss $x_1 \geq 3 \cdot x_2$ gelten, und der Energieverbrauch $E(x_1,x_2)$ soll exakt $E_{\max} = 480$ Wh nicht überschreiten.


a) Klassifiziere für dieses Problem jede der folgenden Größen als \textbf{Variable}, \textbf{Parameter}, \textbf{Zielfunktion}, \textbf{Gleichungs- oder Ungleichungs-Nebenbedingung}: $x_1$, $x_2$, $T$, $w$, $E(x_1,x_2)$, „$x_1 \geq 3x_2$", „$E(x_1,x_2) \leq 480$", „$50 \leq x_1 \leq 200$".


b) Schreibe das vollständige Optimierungsproblem in der Standardform aus dem Kapitel auf (mit $\min f(\vec x)$, $g_j(\vec x) \leq 0$, Boxbedingungen).


c) Begründe, warum „$x_1 \geq 3x_2$" \textbf{keine} Gleichungs-Nebenbedingung sein darf, obwohl man sie auch als „$x_1 - 3x_2 = z$" mit Slack $z \geq 0$ schreiben könnte. Welche Folge hätte die strikte Gleichheit für den zulässigen Bereich $X_z$?


\textbf{Lösung:}


a)

\begin{itemize}
  \item $x_1, x_2$ → \textbf{Variablen} (gezielt variiert)
  \item $T, w$ → \textbf{Parameter} (Umgebung, nicht steuerbar)
  \item $E(x_1,x_2)$ → \textbf{Zielfunktion} $f(\vec x)$
  \item „$x_1 \geq 3x_2$" → \textbf{Ungleichungs-NB}
  \item „$E(x_1,x_2) \leq 480$" → \textbf{Ungleichungs-NB}
  \item „$50 \leq x_1 \leq 200$" → \textbf{Boxbedingungen} ($x_i^l \leq x_i \leq x_i^u$), formal eine spezielle Ungleichungs-NB
\end{itemize}

b) Mit $\vec x = (x_1, x_2)$:

\$\$

\textbackslash{}min\textbackslash{} f(\textbackslash{}vec x) = E(x\_1, x\_2)

\$\$

unter

\$\$

g\_1(\textbackslash{}vec x) = 3x\_2 - x\_1 \textbackslash{}leq 0,\textbackslash{}qquad g\_2(\textbackslash{}vec x) = E(x\_1,x\_2) - 480 \textbackslash{}leq 0,

\$\$

\$\$

50 \textbackslash{}leq x\_1 \textbackslash{}leq 200,\textbackslash{}quad 5 \textbackslash{}leq x\_2 \textbackslash{}leq 20.

\$\$

Es gibt keine Gleichungs-NB ($l=0$), $m=2$ Ungleichungs-NBn (zzgl. Boxen).


c) Die Sicherheitsregel verlangt nur \textbf{„mindestens"} das Dreifache der Geschwindigkeit als Höhe. Eine Gleichheit $x_1 = 3x_2$ würde den Suchraum auf eine Linie kollabieren — alle Punkte mit $x_1 > 3x_2$ wären unzulässig, obwohl sie real sicher sind. $X_z$ wäre dann \textbf{eindimensional} statt zweidimensional, und das Optimum, das im Inneren liegen könnte, würde künstlich ausgeschlossen.



\noindent\textcolor{rulegray}{\rule{\linewidth}{0.4pt}}


\paragraph{Aufgabe 2 — Stationäre Punkte und Hesse \textit{(20 Punkte)}}\mbox{}\\[2pt]

Gegeben sei

\[
f(x_1, x_2) = x_1^3 - 3x_1 + x_2^2 - 4x_2 + 7.
\]


a) Bestimme alle stationären Punkte über die notwendige Bedingung.


b) Klassifiziere jeden stationären Punkt anhand der Hessematrix (Min, Max, Sattel).


c) Erkläre kurz, warum die Bedingung $\nabla f = \vec 0$ allein \textbf{nicht ausreicht}, um ein lokales Minimum zu identifizieren — illustriere mit einem Punkt aus (b).


\textbf{Lösung:}


a) Gradient:

\$\$

\textbackslash{}nabla f(\textbackslash{}vec x) = \textbackslash{}begin\{pmatrix\} 3x\_1\textasciicircum{}2 - 3 \textbackslash{}\textbackslash{} 2x\_2 - 4 \textbackslash{}end\{pmatrix\} \textbackslash{}stackrel\{!\}\{=\} \textbackslash{}vec 0.

\$\$

Aus Komponente 1: $x_1^2 = 1 \Rightarrow x_1 \in \{-1, +1\}$.

Aus Komponente 2: $x_2 = 2$.

Stationäre Punkte: $\vec x_A = (1, 2)$ und $\vec x_B = (-1, 2)$.


b) Hessematrix:

\$\$

H(\textbackslash{}vec x) = \textbackslash{}begin\{pmatrix\} 6x\_1 \& 0 \textbackslash{}\textbackslash{} 0 \& 2 \textbackslash{}end\{pmatrix\}.

\$\$

Da $H$ diagonal ist, sind die Eigenwerte direkt $\lambda_1 = 6x_1,\ \lambda_2 = 2$.


\begin{itemize}
  \item $\vec x_A = (1,2)$: EW $= 6, 2$ → beide $> 0$ → \textbf{positiv definit} → \textbf{lokales Minimum}.
  \item $\vec x_B = (-1,2)$: EW $= -6, 2$ → gemischte Vorzeichen → \textbf{indefinit} → \textbf{Sattelpunkt}.
\end{itemize}

c) $\nabla f = \vec 0$ ist \textbf{notwendig}, aber nicht \textbf{ausreichend}: am Sattelpunkt $\vec x_B$ verschwindet der Gradient ebenfalls, obwohl in $x_1$-Richtung $f$ abfällt und in $x_2$-Richtung steigt. Erst die Hesseanalyse trennt Minima, Maxima und Sattelpunkte voneinander.



\noindent\textcolor{rulegray}{\rule{\linewidth}{0.4pt}}


\paragraph{Aufgabe 3 — Gradient Descent durchrechnen \textit{(20 Punkte)}}\mbox{}\\[2pt]

Gegeben:

\[
f(x_1,x_2) = (x_1 - 3)^2 + 4(x_2 + 1)^2,
\]

Startpunkt $\vec x_0 = (0, 0)$, Abbruchschwelle $\|\nabla f\| < 0{,}5$.


a) Berechne $\nabla f$ symbolisch und bestimme den globalen Minimierer analytisch (welcher Algorithmen-Typ ist das laut Kapitel?).


b) Führe \textbf{drei Iterationen} Gradient Descent mit Schrittweite $c = 0{,}1$ aus. Notiere für $i=0,1,2$: $\vec x_i$, $\nabla f(\vec x_i)$, $\|\nabla f(\vec x_i)\|$. Ist nach drei Schritten abgebrochen?


c) Wiederhole Schritt $i=0 \to i=1$ mit $c = 0{,}5$. Was beobachtest du im Vergleich zu (b), und welcher der im Kapitel genannten \textbf{Probleme von Gradient Descent} liegt vor?


\textbf{Lösung:}


a) $\nabla f = \begin{pmatrix} 2(x_1-3) \\ 8(x_2+1) \end{pmatrix}$. Setze $=\vec 0$: $x_1 = 3,\ x_2 = -1$. Globaler Minimierer $\vec x^* = (3,-1)$, $f(\vec x^*) = 0$.

Da rein per Symbolrechnen ohne Iteration gelöst → \textbf{analytisch} und \textbf{exakt} (Klassifikation aus dem Kapitel).


b) Update-Regel: $\vec x_{i+1} = \vec x_i - 0{,}1 \cdot \nabla f(\vec x_i)$.



{\small
\begin{tabularx}{\linewidth}{XXXXXX}
\hline
\rowcolor{tableheadbg}
{\color{white}\textbf{$i$}} & {\color{white}\textbf{$\vec x_i$}} & {\color{white}\textbf{$\nabla f(\vec x_i)$}} & {\color{white}\textbf{\$\textbackslash{}}} & {\color{white}\textbf{\textbackslash{}nabla f\textbackslash{}}} & {\color{white}\textbf{\$}} \\[2pt]
\hline
\endhead
\hline
\endfoot
0 & $(0,\ 0)$ & $(-6,\ 8)$ & $10{,}00$ &  &  \\
1 & $(0{,}6,\ -0{,}8)$ & $(-4{,}8,\ 1{,}6)$ & $5{,}06$ &  &  \\
2 & $(1{,}08,\ -0{,}96)$ & $(-3{,}84,\ 0{,}32)$ & $3{,}85$ &  &  \\
3 & $(1{,}464,\ -0{,}992)$ & — & — &  &  \\
\end{tabularx}
}


Rechnung Schritt 1: $\vec x_1 = (0,0) - 0{,}1 \cdot (-6, 8) = (0{,}6, -0{,}8)$.

Schritt 2: $\nabla f(0{,}6,-0{,}8) = (2(0{,}6-3),\ 8(-0{,}8+1)) = (-4{,}8,\ 1{,}6)$; $\vec x_2 = (0{,}6,-0{,}8) - 0{,}1(-4{,}8,\ 1{,}6) = (1{,}08,\ -0{,}96)$.

Schritt 3: $\nabla f(1{,}08,-0{,}96) = (-3{,}84,\ 0{,}32)$; $\vec x_3 = (1{,}08,-0{,}96) - 0{,}1(-3{,}84,\ 0{,}32) = (1{,}464,\ -0{,}992)$.


Da $\|\nabla f(\vec x_2)\| = 3{,}85 \not< 0{,}5$, wird \textbf{nicht abgebrochen}.


c) Mit $c = 0{,}5$: $\vec x_1 = (0,0) - 0{,}5\cdot(-6,8) = (3, -4)$.

$\nabla f(3,-4) = (0,\ -24)$, $\|\nabla f\| = 24$.

In $x_1$ wurde das Optimum perfekt getroffen, in $x_2$ aber \textbf{weit überschossen} (von 0 auf -4, Optimum bei -1). Im nächsten Schritt würde $x_2$ wieder zurückspringen → \textbf{Oszillation} wegen zu großem $c$ in der $x_2$-Richtung (dort ist die Krümmung 4× so hoch wie in $x_1$). Genau das im Kapitel genannte Problem „falsches $c$ → Oszillation".



\noindent\textcolor{rulegray}{\rule{\linewidth}{0.4pt}}


\paragraph{Aufgabe 4 — Eigenschaften zweier Verfahren vergleichen \textit{(15 Punkte)}}\mbox{}\\[2pt]

Du sollst eine Bauteil-Anordnung optimieren. Variante A: alle Designparameter sind \textbf{kontinuierlich}, $f$ ist glatt und differenzierbar, vermutlich \textbf{unimodal}. Variante B: die Parameter sind \textbf{diskret/kombinatorisch} (Reihenfolge auf einer Platine, Permutation), $f$ ist \textbf{multimodal} und nicht ableitbar; jede Auswertung ist zudem leicht \textbf{verrauscht} (Messrauschen am Prüfstand).


a) Ordne \textbf{Gradient Descent} und einen \textbf{Evolutionären Algorithmus (EA)} anhand der Tabelle „Eigenschaften (Algorithmen vs. Probleme)" aus dem Kapitel jeweils in mindestens fünf Kategorien ein (lokal/global, Differenzierbarkeit, Linearität, Genauigkeit, Zufallseinfluss, Variablenart).


b) Welcher Algorithmus passt zu Variante A, welcher zu B? Begründe \textbf{pro Kategorie} mit der Tabelle.


c) Was passiert, wenn man Gradient Descent trotzdem auf Variante B loslässt? Gib mindestens zwei konkrete Mismatches an.


\textbf{Lösung:}


a)


{\small
\begin{tabularx}{\linewidth}{XXX}
\hline
\rowcolor{tableheadbg}
{\color{white}\textbf{Eigenschaft}} & {\color{white}\textbf{Gradient Descent}} & {\color{white}\textbf{EA}} \\[2pt]
\hline
\endhead
\hline
\endfoot
lokal/global & lokal & global \\
Differenzierbarkeit & gradientenbasiert & gradientenfrei \\
Linearität & für nicht-lineare $f$ einsetzbar (NB: linear-frei) & nicht-linear \\
Lösungsansatz & numerisch & numerisch \\
Genauigkeit & approximativ & approximativ \\
Zufallseinfluss & deterministisch & stochastisch \\
Variablen & kontinuierlich & dis./kom./kon./gem. \\
\end{tabularx}
}


b)

\begin{itemize}
  \item \textbf{Variante A → Gradient Descent}:
\begin{itemize}
  \item unimodal ↔ lokal reicht aus,
  \item differenzierbar ↔ Gradient verfügbar,
  \item kontinuierlich ↔ passt zur Variablenart,
  \item kein Rauschen vorausgesetzt ↔ deterministisch ok.
\end{itemize}
  \item \textbf{Variante B → EA}:
\begin{itemize}
  \item multimodal ↔ globaler Algorithmus nötig,
  \item nicht differenzierbar ↔ gradientenfrei,
  \item kombinatorisch ↔ EA kann Permutationen/Strings,
  \item stochastisches Problem ↔ stochastischer Algorithmus robuster (z. B. durch Populationsmittel).
\end{itemize}
\end{itemize}

c) Mismatches GD ↔ Variante B:

\begin{enumerate}
  \item \textbf{Kein Gradient definiert} auf Permutationen → Update-Regel $\vec x_{i+1} = \vec x_i - c\nabla f$ nicht anwendbar.
  \item \textbf{Nur lokal} → bei Multimodalität bleibt GD im erstbesten lokalen Min hängen.
  \item \textbf{Deterministisch + Rauschen}: GD reagiert direkt auf jede verrauschte Gradientenschätzung → Konvergenzprobleme.
\end{enumerate}


\noindent\textcolor{rulegray}{\rule{\linewidth}{0.4pt}}


\paragraph{Aufgabe 5 — Hesse-Klassifikation am 3-Hump-Beispiel \textit{(15 Punkte)}}\mbox{}\\[2pt]

Gegeben ist die im Kapitel verwendete 3-Hump-Funktion

\[
f(x_1,x_2) = 2x_1^2 - 1{,}05\, x_1^4 + \tfrac{1}{6}x_1^6 + x_1 x_2 + x_2^2.
\]

Der Gradient steht im Kapitel:

$\nabla f = \big(4x_1 - 4{,}2\, x_1^3 + x_1^5 + x_2,\ \ x_1 + 2x_2\big)^\top.$


a) Verifiziere, dass $\vec x^* = (0,0)$ ein stationärer Punkt ist.


b) Berechne die Hessematrix $H(\vec x)$ allgemein und werte sie an $\vec x^*$ aus.


c) Bestimme die Eigenwerte von $H(0,0)$ und klassifiziere $\vec x^*$. Welche der vier Hesse-Fälle aus der Tabelle des Kapitels liegt vor?


\textbf{Lösung:}


a) $\nabla f(0,0) = (4\cdot 0 - 0 + 0 + 0,\ 0 + 0) = (0,0)$ ✓ → SP.


b) Zweite Ableitungen:

\begin{itemize}
  \item $\partial^2 f/\partial x_1^2 = 4 - 12{,}6\, x_1^2 + 5x_1^4$
  \item $\partial^2 f/\partial x_2^2 = 2$
  \item $\partial^2 f/\partial x_1 \partial x_2 = 1$
\end{itemize}

\$\$

H(\textbackslash{}vec x) = \textbackslash{}begin\{pmatrix\} 4 - 12\{,\}6 x\_1\textasciicircum{}2 + 5 x\_1\textasciicircum{}4 \& 1 \textbackslash{}\textbackslash{} 1 \& 2 \textbackslash{}end\{pmatrix\},\textbackslash{}qquad

H(0,0) = \textbackslash{}begin\{pmatrix\} 4 \& 1 \textbackslash{}\textbackslash{} 1 \& 2 \textbackslash{}end\{pmatrix\}.

\$\$


c) Charakteristisches Polynom:

$\det(H - \lambda I) = (4-\lambda)(2-\lambda) - 1 = \lambda^2 - 6\lambda + 7 = 0.$

$\lambda_{1,2} = 3 \pm \sqrt{9-7} = 3 \pm \sqrt{2} \approx 4{,}414;\ 1{,}586.$


Beide Eigenwerte $> 0$ → \textbf{positiv definit} → laut Tabelle des Kapitels: \textbf{lokales Minimum} an $(0,0)$.



\noindent\textcolor{rulegray}{\rule{\linewidth}{0.4pt}}


\paragraph{Aufgabe 6 — Pseudocode-Fehler im Gradient Descent \textit{(15 Punkte)}}\mbox{}\\[2pt]

Ein Kommilitone hat folgenden Pseudocode für Gradient Descent abgegeben:


\begin{lstlisting}
x = x0
c = 0.1
for i = 1 .. 1000:
    g = gradient(f, x)
    x = x + c * g                # (1)
    if norm(g) > 1e-3:           # (2)
        break
return x
\end{lstlisting}


a) Identifiziere \textbf{mindestens zwei} inhaltliche Fehler bzgl. des im Kapitel beschriebenen Algorithmus. Gib jeweils die Zeile und die Korrektur an.


b) Erkläre an Zeile (1), warum der falsche Operator in der Praxis nicht „nur langsam", sondern \textbf{systematisch falsch} konvergiert (was tut der Algorithmus damit eigentlich?).


c) Skizziere kurz, wie du den korrigierten Algorithmus auf die Funktion $f(x) = (x-2)^2$ mit Startpunkt $x_0 = 5$, $c = 0{,}25$ anwenden würdest — gib die ersten zwei Iterationen an.


\textbf{Lösung:}


a)

\begin{itemize}
  \item \textbf{Zeile (1)} \texttt{x = x + c*g} ist falsch. Im Kapitel: $\vec x_{i+1} = \vec x_i - c\cdot\nabla f(\vec x_i)$. Korrektur: \texttt{x = x - c * g}.
  \item \textbf{Zeile (2)} Bedingung ist invertiert: Bei Gradient Descent soll abgebrochen werden, wenn $\|\nabla f\| \approx 0$, also \textbf{klein}. Der Code bricht ab, sobald der Gradient \textbf{groß} ist. Korrektur: \texttt{if norm(g) < 1e-3: break}.
  \item (Optional) Es fehlt eine Schleife/Rückgabe nach dem Break — okay, vorhanden, aber: der Test sollte \textbf{vor} dem Update stehen, sonst macht man immer einen Schritt zu viel (entspricht Schritt 3 im Kapitel-Pseudocode).
\end{itemize}

b) \texttt{x = x + c*g} bewegt sich in Richtung des \textbf{stärksten Anstiegs} (das ist genau die Definition von $\nabla f$). Damit führt der Algorithmus \textbf{Gradient Ascent} aus und maximiert $f$ statt zu minimieren. Bei einer konvexen Funktion läuft $\vec x$ also \textbf{systematisch ins Unendliche} weg vom Minimum, nicht nur langsam.


c) $f(x) = (x-2)^2$, $f'(x) = 2(x-2)$.

\begin{itemize}
  \item $i=0$: $x_0 = 5$, $g_0 = 2(5-2) = 6$, $x_1 = 5 - 0{,}25 \cdot 6 = 3{,}5$.
  \item $i=1$: $g_1 = 2(3{,}5-2) = 3$, $x_2 = 3{,}5 - 0{,}25\cdot 3 = 2{,}75$.
\end{itemize}

Konvergiert monoton gegen $x^* = 2$.





% ──────────────────────────────────────────────────────────────

\end{document}
